\section{Automatic column scales and zero baselines}
\label{sec:column-scale}
Column width controls available space; the numerical scale controls where values
appear within that space. Table and column keys \key{scale=fixed|column} and
\key{include zero=true|false} control ranges independently of width.
Fixed is the default and retains min/max in drawing options. Column mode scans
every row before layout and overrides those limits. Column settings override
table settings; text columns ignore scaling. Manual and CSV rows use the same
rules, with consistent ranges across pages. CSV selection defaults to full-data
ranges; \key{scale scope=visible} selects the displayed rows instead.
\begin{demo}{Shared ranges with per-column zero settings}
\begin{gradtable}[width=\linewidth,scale=column,include zero=false,header align=c]
  \gradcolumn{name}{Method}
  \gradcolumn[type=bar,include zero=true]{score}{Score}
  \gradcolumn[type=bar]{time}{Time}
  \gradrow{A,70,70}
  \gradrow{B,90,90}
  \gradrow{Pending,NA,NA}
\end{gradtable}
\end{demo}
Score uses [0,90]; Time uses [70,90].
\subsection{Including zero or retaining a fixed range}
Missing values are ignored; zero remains an observation and is included by
default. Without zero, positive-only and negative-only ranges are valid; a bar's
zero baseline is clipped to the nearest range border. Keep zero when lengths
must describe magnitudes from zero. Percent labels still require \key{min=0} and max>0.
Use \key{scale=fixed} with the same limits when comparing different datasets. Targets,
references and error magnitudes are not observations used for scaling; choose a
fixed range when those annotations need additional space.
\subsection{Missing, constant and compound data}
All-missing or header-only columns use [0,1]; zero-inclusive all-zero columns
also use [0,1]. A constant column without zero uses [value-0.5,value+0.5].
Number, bar, bullet and heat columns use scalar extrema. Pair and interval
columns use endpoints and point estimates; spark columns share every observation
and select a fixed range; stacks use signed cumulative totals with zero baselines.

\subsection{Shared named scales across tables}
\key{\gradscale{name}{min=lo,max=hi}} declares explicit limits with min<max.
Scale group references work in standalone graphics or table columns. Definitions
follow TeX grouping; redefinition replaces the current scope's value. Groups
override drawing min/max options and do not depend on table order.
\begin{demo}{A graphic and table sharing an accuracy scale}
\gradscale{accuracy}{min=0,max=100}
\gradbar[scale group=accuracy,theme=violet]{72}

\begin{gradtable}[style=academic,theme=violet,width=\linewidth,header align=c]
  \gradcolumn{name}{Method}
  \gradcolumn[type=bar,scale group=accuracy]{score}{Score}
  \gradrow{A,72}\gradrow{B,90}
\end{gradtable}
\end{demo}
Table/column groups require \key{scale=fixed}; \key{scale=column} is a conflict. Columns
inherit table groups and can clear them with an empty scale group. Ordinary
sparks select fixed group ranges; win-loss retains sign categories. Unknown
groups are errors. This release does not automatically collect future tables.
\subsection{Padding, nice bounds and percent denominators}
Scale padding and scale nice apply only to automatic column ranges and default
to 0 and false. Padding is a percentage of the original span at each nonzero
endpoint, retaining zero anchors. Nice bounds round outward to a 1/2/5 step,
targeting about five intervals. Column settings override table settings; fixed
or named limits are unchanged.
\begin{demo}{Expanded drawing limits with the original percent basis}
\begin{gradtable}[style=compact,theme=forest,width=\linewidth,
  scale=column,scale padding=10,scale nice=true,header align=c]
  \gradcolumn{name}{Method}
  \gradcolumn[type=bar,aggregate=mean,
    options={value format=percent}]{score}{Score}
  \gradrow{A,63}\gradrow{B,87}
  \gradsummary{Mean}
\end{gradtable}
\end{demo}
The range starts at [0,87], expands to [0,95.7] and rounds to [0,100]. Percent
labels still use 87, so the old maximum stays 100\%. Drawing key percent base
can explicitly select another positive denominator. Without adjustment the
existing max-based behavior remains. Graphic percent formatting requires \key{min=0}.
Ranks and aggregates use raw values; heat centers must remain inside the range.
